\section{Ground-truth construction}
\label{sec:groundtruth}

The validity of everything that follows depends on this section, so we describe the
corpus in more detail than is conventional. Two design commitments govern it. First,
no label originates from human judgement about whether a text looks machine-written;
every label derives from platform attribution or from chronology. Second, positives and
negatives are drawn from the same repositories wherever possible, so that differences
between them are not differences between projects.

\subsection{Positive set: declared agent authorship}
\label{sec:groundtruth:positives}

Positives come from the AIDev dataset \citep{li2025aidev}, which collects pull requests
opened by agentic coding platforms operating under their own accounts. We pin the
dataset revision used, because AIDev is a living resource: the paper introducing it
reports \citednum{456{,}000} pull requests, while the revision we analyse contains \AidevPRs{}
across \AidevRepos{} repositories. All counts below are measured from the pinned
revision rather than quoted from the literature.

Three properties of the positive set constrain the design, and we report them
prominently because they bound every claim in this paper.

\paragraph{The observation window is seven months.} All agent-authored pull requests in
the dataset fall between \AidevWindowStart{} and \AidevWindowEnd{}. Our findings
therefore concern the agent platforms and underlying model versions current in the first
half of 2025. We return to this in Section~\ref{sec:threats}.

\paragraph{The agent mix is severely imbalanced.} Table~\ref{tab:agents} gives the
composition. A single agent accounts for \AidevTopAgentShare{} of all pull requests, and
median body length differs by roughly an order of magnitude across agents. Any metric
computed by pooling agents is therefore dominated by one of them. We report all
detector results per agent, and present pooled figures only under explicit equal
allocation.

\begin{table}
\caption{Composition of the positive set at the pinned AIDev revision. Median body
length varies by roughly an order of magnitude across agents, so pooled metrics are
dominated by the largest contributor.}
\label{tab:agents}
\begin{tabular}{lrr}
\hline\noalign{\smallskip}
Agent & Pull requests & Median body (chars) \\
\noalign{\smallskip}\hline\noalign{\smallskip}
OpenAI Codex & \AgentOpenAICodexPRs{} & \AgentOpenAICodexBodyMedian{} \\
GitHub Copilot & \AgentCopilotPRs{} & \AgentCopilotBodyMedian{} \\
Cursor & \AgentCursorPRs{} & \AgentCursorBodyMedian{} \\
Devin & \AgentDevinPRs{} & \AgentDevinBodyMedian{} \\
Claude Code & \AgentClaudeCodePRs{} & \AgentClaudeCodeBodyMedian{} \\
\noalign{\smallskip}\hline
\end{tabular}
\end{table}

\paragraph{The texts are short.} The median pull request body is \AidevBodyMedian{}
characters and the median commit message \AidevCommitMsgMedian{}. This is not incidental
to the study; it is arguably the central fact about the domain. Detectors that report
near-perfect separation on essay-length text are being asked here to decide on a single
line of English.

\subsubsection{Artifact genres}
\label{sec:groundtruth:genres}

We treat five genres separately, because they differ in length, register and
conventionality, and because the prevalence literature scores different subsets of them:

\begin{enumerate}
  \item \emph{Pull request descriptions} --- prose, moderate length, often templated.
  \item \emph{Commit messages} --- the shortest genre, heavily conventionalised.
    \texttt{Co-authored-by} and similar trailer lines are stripped into a separate field
    and never passed to a detector; recovering an explicit declaration is not detection.
  \item \emph{Code diffs} --- the patch hunks, scored per file, with language taken from
    repository metadata.
  \item \emph{Issue bodies} --- prose with a different register from pull requests.
  \item \emph{Pull request comments} --- restricted to human accounts. This restriction
    is severe and necessary: \AidevCommentBotShare{} of comments on agent-authored pull
    requests are bot-authored, so an unfiltered comment genre would measure
    continuous-integration bot output rather than agent output.
\end{enumerate}

Allocation across the resulting agent-by-genre cells is equal rather than proportional,
targeting 1{,}000 texts per cell with a floor of 300. Equal allocation sacrifices
precision on the most abundant agent to obtain usable precision on the least abundant,
which is the correct trade when the question is whether detector behaviour generalises
across agents rather than what its average is. The sampled positive set contains
\PositiveTexts{} texts across \PositiveCells{} cells.

\paragraph{What equal allocation costs.} Because agents differ so much in verbosity,
reweighting them equally changes the corpus itself: the median pull request body rises
from \PrBodyMedianPopulation{} characters in the population to
\PrBodyMedianSampled{} in the sample. The sampled corpus is therefore deliberately
\emph{not} representative of the population of agent contributions, and no prevalence
claim is made from it --- it exists to estimate per-stratum detector error, which
Section~\ref{sec:correction} then applies to representative samples. Reporting this
explicitly matters because the same reweighting, left undisclosed, would make detectors
look better than they are: longer texts are easier to classify.

\paragraph{One genre cannot support per-agent analysis.} Issue bodies are linked to
agent authorship only through AIDev's issue--pull-request link table, and the resulting
cells are thin: three of the five agents fall below the floor of 300, two of them below
200. Issue bodies are therefore analysed pooled across agents, and per-agent claims are
restricted to the other four genres. This is a limitation of the available attribution
data rather than of the design, and we report it rather than quietly padding the cells
from a proportional draw.

\subsection{Negative sets}
\label{sec:groundtruth:negatives}

No single negative set is simultaneously pure and representative, so we construct three
with different and complementary weaknesses.

\subsubsection{N1: pre-LLM artifacts from the same repositories}
\label{sec:groundtruth:n1}

The primary negative set consists of artifacts written in the positive set's own
repositories before \LlmCutoff{}, the public release of ChatGPT. Its purity follows from
chronology, and every detection in it is a false positive by construction.

\paragraph{Dating the text, not the artifact.} Filtering on an artifact's creation date
is insufficient, because bodies can be edited and issue threads accumulate comments
indefinitely. While profiling repository data we encountered an issue created in 2013
carrying a comment written in 2026 that opens with an explicit agent-generated
verification report. Selecting on creation date alone would have admitted that comment
to the pre-LLM negative set. We therefore require both the creation and last-modification
timestamps to precede the cutoff, which guarantees the text has not been touched since,
at the cost of biasing the set toward artifacts that attracted no subsequent activity.
Section~\ref{sec:threats} treats this quiescence bias.

For commit messages and diffs the problem does not arise in the same form: we read them
from cloned repository history rather than from the platform API, where they are dated
by committer timestamp and cannot be altered without rewriting history.

\paragraph{Repository eligibility.} Not every repository in the positive set has a
pre-LLM history. Of the \RepoFrame{} curated repositories, \NOneResolved{} still resolve
on GitHub --- \NOneUnresolvable{} have been deleted, renamed or withdrawn --- and of
those, \NOneEligible{} (\NOneEligiblePct{}) were created before \LlmCutoff{}. These form
the N1 frame. The oldest dates from \NOneOldestRepo{}, so the frame spans the better part
of two decades of development practice.

Repository names are not stable identifiers: a project can be renamed and its name
reused by an unrelated project. We therefore verify that each name still resolves to the
GitHub numeric identifier that AIDev recorded, and exclude any repository where it does
not, independently of age.

The excluded \NOneExcluded{} repositories are not a random subset --- they are younger, and younger projects plausibly differ in their
propensity to adopt coding agents. Losing nearly two-fifths of the frame is a real cost
to statistical power in the matched analysis, and we report a generalisation check on
the excluded repositories using positives alone to test whether the restriction changes
the conclusions.

\subsubsection{N2: contemporaneous human contributions}
\label{sec:groundtruth:n2}

N1 is pure but temporally separated from the positives, which confounds authorship with
period: coding conventions, tooling and project norms all changed between 2022 and 2025.
The second negative set removes the temporal confound at the cost of purity. AIDev
includes \AidevHumanPRs{} human-authored pull requests from \AidevHumanRepos{}
repositories and \AidevHumanUsers{} distinct users, drawn from the same period as the
positives, with a median body of \AidevHumanBodyMedian{} characters.

These are human-\emph{attributed}, not verified human-\emph{written}: a developer may
have drafted the text with a chatbot and pasted it in, leaving no trace. Error rates on
N2 are therefore reported as \emph{apparent} rather than true false positive rates. The
gap between the N1 and N2 rates is itself informative, since under stated assumptions it
bounds the extent of undeclared assisted authorship in the human population --- the
stratum this study otherwise cannot observe.

\subsubsection{N3: pre-LLM automation}
\label{sec:groundtruth:n3}

The third set isolates the confounder. It comprises pre-cutoff artifacts from known
automation accounts --- dependency managers, release tooling, continuous integration
reporters --- in the same repositories. This text is machine-generated and definitively
not LLM-generated. The detection rate on N3 estimates how much reported ``AI
prevalence'' is decade-old automation being relabelled.

\subsection{Matching and stratification}
\label{sec:groundtruth:matching}

Because detector scores are length-sensitive and the agents differ substantially in
length, an unmatched comparison would partly measure length rather than authorship.
Negatives are therefore sampled to match the positive length distribution within genre
by decile. We report unmatched results alongside matched ones, since the difference
between them estimates the magnitude of the length confound rather than merely
controlling for it.

Stratification variables are programming language, text length, artifact genre, agent,
and project scale. For the subset of repositories also covered by the CrOSSD panel
\RESULT{cite Dam, Klausner, Neumaier WWW'23} we additionally use its health and
criticality metrics. Following \citet{liang2023biased}, we require a proxy for
contributors' English proficiency, but author-name and author-location proxies are both
unreliable and ethically objectionable. We instead use repository-level linguistic
proxies derived from documentation text, which attribute a property to the project
rather than to an individual.

\subsection{Corpus release}
\label{sec:groundtruth:release}

The frozen corpus is released as a specification --- artifact identifiers, selection
predicates, seeds and per-cell counts --- together with the code that rebuilds it, so
that the corpus can be reconstructed without redistributing repository content beyond
what its licence permits. Every count reported in this paper is emitted directly by the
analysis scripts into the manuscript source, so no figure in the prose can drift from the
data that produced it.
